\chapter{Introduction}
\label{intro}
{\itshape Real-time video segmentation requires balancing mathematical accuracy with strict latency budgets. This work introduces a comparative framework of classical and quantum-inspired 5D clustering algorithms implemented on the Graphics Processing Unit (GPU). By analyzing the computational limits of CPU (Central Processing Unit) architectures and physical Noisy Intermediate-Scale Quantum (NISQ) hardware, the engineering objectives and concrete software contributions of this implementation are established.}

\section{Context and Motivation}
Image segmentation is a fundamental computer vision task that partitions an image into meaningful, uniformly connected regions \cite{gonzalez2009digital}. It is a key preprocessing step for high-level computer vision applications, such as lane-tracking in autonomous vehicles \cite{janai2020computer}, object detection \cite{redmon2016you}, and medical imaging overlays \cite{pham2000current}. While simple color-based segmentation is common, modern real-time applications require spatial consistency. To achieve this, each pixel is mapped to a 5D feature manifold:
\begin{equation}
\mathbf{x} = [B, G, R, x, y]^T
\end{equation}
By clustering pixels within this combined color-spatial manifold, the algorithm naturally groups pixels that are both chromatically similar and geometrically close, preventing pixel-level noise \cite{achanta2012slic}.

Real-time video processing places strict constraints on execution latency. In computer vision, standard real-time throughput is governed by the human-facing display refresh rate of 60 frames per second (FPS) ($16.67\text{ ms}$ per frame) \cite{DHANANI20135}. Yet, for high-speed interactive systems, including driving simulators \cite{cruden2019simulator}, virtual reality headsets \cite{jerald2009scene}, and by extension robotic manipulation loops and autonomous vehicle obstacle avoidance, the standard target is high-speed real-time processing operating at $120\text{ FPS}$ or more ($8.33\text{ ms}$ or less per frame) to prevent feedback loop delays and tracking oscillation. Running heavy deep-learning segmentation models on embedded edge devices is often impossible under these extreme limits due to massive power consumption and high latency \cite{paszke2016enet,slimani2019k}. Therefore, highly optimized, lightweight, GPU-accelerated classical clustering is a practical alternative for edge devices \cite{mucke2019hardware}.

Simultaneously, quantum-inspired computing offers a practical model for simulating quantum mechanics on classical hardware. True quantum computers are still in the NISQ era \cite{preskill2018quantum}, characterized by small physical qubit registers, lack of error correction, and short coherence times ($T_1, T_2$) \cite{bharti2022noisy}. Physical hardware is highly sensitive to environmental noise, and transferring live video frames to a quantum cloud service introduces massive network latency, rendering genuine quantum hardware unusable for real-time video \cite{fraunhofer2025quantum}.

To bypass these hardware constraints while leveraging quantum mathematical behaviors, quantum-inspired computing is utilized. This approach translates quantum mechanics such as state-overlap fidelity via the SWAP test \cite{buhrman2001quantum} (which measures $|\langle\psi\vert\phi\rangle|^2$) into deterministic equations that are parallelized directly inside custom GPU kernels. This allows for a direct comparison of the mathematical properties of classical and quantum clustering in real-time on standard local hardware.

\section{Problem Statement}
Designing a real-time, comparative 5D clustering framework introduces three major technical challenges:
\begin{enumerate}
    \item \textbf{Computational Complexity:} Evaluating 5D similarity for a standard Video Graphics Array (VGA) frame ($640 \times 480$ pixels, $N = 307,200$ pixels) across $K$ clusters requires millions of floating-point calculations per frame. Doing this sequentially on a CPU is too slow for real-time feeds, resulting in massive lag.
    \item \textbf{NISQ Hardware Limitations:} While quantum algorithms can evaluate state overlaps efficiently, physical Quantum Processing Units (QPUs) cannot operate within the sub-millisecond execution windows required for video pipelines due to hardware noise \cite{bharti2022noisy} and slow communication interfaces \cite{fraunhofer2025quantum}.
    \item \textbf{Lack of a Unified Benchmark:} To the best of the author's knowledge, there is no standard, open-source tool that allows for a direct, fair, and real-time comparison between classical Euclidean distance \cite{lloyd1982least} and quantum-inspired phase-space metrics under identical initialization and hardware constraints.
\end{enumerate}

\section{Thesis Objectives}
This thesis develops and validates a high-throughput, GPU-accelerated C++/CUDA (Compute Unified Device Architecture) framework for real-time 5D picture segmentation. The specific objectives are to:
\begin{itemize}
    \item \textbf{Build an Asynchronous Pipeline:} Implement a multi-threaded execution model in C++20 that decouples raw video frame acquisition (OpenCV \cite{Bradski2000}) from the GPU computational backend to target high-speed processing of at least $120\text{ FPS}$ ($8.33\text{ ms}$ processing budget per frame) while keeping the Graphical User Interface (GUI) highly responsive.
    \item \textbf{Implement Interchangeable GPU Engines:} Develop modular, parallel computational backends for both classical Euclidean distance and quantum-inspired fidelity metrics within a clean Strategy design pattern \cite{gamma1995design}.
    \item \textbf{Maximize VRAM (Video Random Access Memory) Efficiency:} Design custom CUDA kernels that utilize high-speed on-chip shared memory atomics to reduce global memory contention during centroid updates \cite{CudaGuide}.
    \item \textbf{Design a Live Control Dashboard:} Build an interactive control panel using Dear ImGui \cite{ImGui} to enable real-time hyperparameter adjustments (number of clusters $K$, learning interval, and stride $S$), engine hot-swapping, and the monitoring of core GPU execution speed and overall pipeline throughput.
    \item \textbf{Implement On-Demand Quality Diagnostics:} Develop an asynchronous diagnostic utility that captures a snapshot frame from the active video feed to evaluate segmentation quality (Within-Cluster Sum of Squares, or WCSS \cite{lloyd1982least, krzanowski1988criterion}, the Davies-Bouldin ($\mathcal{D}\mathcal{B}$) index \cite{davies1979cluster,xiao2017davies}, and a Monte Carlo-approximated Silhouette score \cite{shutaywi2021silhouette, rousseeuw1987silhouettes}) alongside convergence iteration counts and stage latencies on a background thread, presenting the results in a User Interface (UI) overlay without interrupting the main video pipeline.
\end{itemize}

\section{Original Contributions}
This paper presents a fully verified, open-source C++/CUDA software package. The original contributions of this work include:
\begin{enumerate}
    \item \textbf{An Optimized Quantum-Inspired Cosine Kernel:} The emulation of the quantum SWAP test circuit \cite{buhrman2001quantum} on the GPU by mapping pixel coordinates to phase states and calculating overlaps using direct hardware-accelerated cosine intrinsics (\texttt{\_\_cosf()}), leveraging Special Function Unit (SFU) hardware acceleration \cite{CudaGuide} to achieve high-throughput single-instruction SASS (Streaming Assembly) execution.
    \item \textbf{A Modular Strategy Engine Architecture:} The design of a fully decoupled C++ execution framework based on the Strategy design pattern \cite{gamma1995design}, allowing for seamless, real-time hot-swapping between the classical Euclidean and quantum-inspired GPU engines without resetting the video stream or losing state.
    \item \textbf{An Asynchronous Still-Frame Diagnostic Engine:} The development of an on-demand diagnostic worker thread that captures a snapshot still-frame from the active feed to calculate heavy validation metrics (WCSS \cite{lloyd1982least,krzanowski1988criterion}, the Davies-Bouldin ($\mathcal{D}\mathcal{B}$) index \cite{davies1979cluster,xiao2017davies}, and a Monte Carlo-approximated Silhouette score \cite{rousseeuw1987silhouettes, shutaywi2021silhouette}) alongside convergence iteration counts and stage latencies, presenting this complete diagnostic report in a dedicated UI overlay without lagging the live video feed.
    \item \textbf{A Parallel Local-to-Global Atomic Reduction Loop:} The implementation of a two-tier aggregation strategy that accumulates sub-totals in fast shared memory before updating global VRAM, reducing global memory traffic by a factor of $256\times$ \cite{CudaGuide}.
    \item \textbf{A Massively Parallel Classical GPU Kernel:} The implementation of an optimized classical Euclidean 5D K-Means clustering kernel utilizing shared-memory caching for cluster centroids and cache-optimized global memory accesses to maximize GPU execution throughput \cite{CudaGuide}.
    \item \textbf{A Custom GPU-Native Preprocessor:} The development of the \texttt{StridedData\allowbreak Preprocessor}, which normalizes Blue-Green-Red (BGR) and spatial pixel coordinates and formats them into a contiguous Array of Structures (AoS) on the GPU, maximizing $L_1$ cache hit rates \cite{CudaGuide}.
    \item \textbf{An Interactive Control Dashboard:} The creation of a professional-grade Dear ImGui \cite{ImGui} control panel positioned alongside the active video feed to display live performance diagnostics (FPS, algorithm execution speed, and camera pipeline retrieval speed) while enabling users to dynamically hot-swap engines, tune hyperparameters ($K$ clusters, learning interval, and stride $S$), toggle centroid visualizations, and reset active cluster states in real-time.
    \item \textbf{A Rigorous Testing Suite and Fixed-Seed Protocol:} The construction of a comprehensive automated unit-testing suite in GoogleTest (GTest) \cite{googletest} to guarantee code parity, empty cluster resilience, and floating-point stability using a global reproducibility seed (\texttt{STABLE\_RANDOM\_SEED = 42}).
\end{enumerate}

\section{Thesis Structure}
The rest of this thesis is organized as follows:
\begin{itemize}
    \item \textbf{Chapter 2 (State of the Art \& Theory):} Reviews classical K-Means, multi-modal scaling, CUDA Single Instruction, Multiple Threads (SIMT) execution models, NISQ physics limits, and the quantum SWAP test formulation.
    \item \textbf{Chapter 3 (System Architecture \& Design):} Illustrates the asynchronous producer-consumer model, C++ design patterns (Strategy, Facade, Pimpl), and the 5D data engineering pipeline.
    \item \textbf{Chapter 4 (Implementation and Optimization):} Details NVCC (NVIDIA CUDA Compiler) compiler optimizations, shared memory atomics, hardware intrinsics, and page-locked Direct Memory Access (DMA) staging.
    \item \textbf{Chapter 5 (Verification and Testing):} Presents the automated GTest suite, boundary condition assertions, CPU-GPU parity testing, and initial performance telemetry.
    \item \textbf{Chapter 6 (Results and Comparative Analysis):} Compares the classical and quantum-inspired engines in terms of convergence speed, image segmentation detail, and telemetry metrics.
    \item \textbf{Chapter 7 (Conclusions and Future Work):} Summarizes the engineering findings and outlines future optimization routes on embedded systems and physical QPUs.
\end{itemize}

\section{Declaration of Generative AI Usage}
During the preparation of this work, the author used Google Gemini in order to improve language readability, assist in structuring the chapter outlines, and refine the articulation of technical concepts. Regarding the software implementation, the technology assisted in maintaining code quality and offering guidance for algorithmic performance optimization. After using this tool, the author reviewed and edited the content as needed, and takes full responsibility for both the text and the final implementation of the thesis.
